
\documentclass[a4paper,12pt]{article}

\usepackage{amsthm}
\usepackage{natbib}
\usepackage{amsmath}

\newtheorem{theorem}{Theorem}[section]

\begin{document}
	
	\title{My First LaTeX Document}
	\author{Manasa Paulson}
	\date{\today}
	
	\maketitle
	
	\tableofcontents
	
	\newpage
	
	\section{Introduction}\label{intro}
	
	Words are separated by one or more spaces. Paragraphs are
	separated by one or more blank lines. The output is not affected
	by adding extra spaces or extra blank lines to the input file.
	
	{\large $\beta = a^2 + b\times c + 2$}
	
	\subsection{Background}
	
	In Section \ref{intro}, double quotes are typed like this:
	``quoted text''.
	
	Single quotes are typed like this: `single-quoted text'.
	
	\textbf{Long dashes} are typed as three dash characters---like this.
	
	Italic text is typed like this: \textit{This is italic text}.
	
	Bold text is typed like this: \textbf{This is bold text}.
	
	\subsection{A Numbered List}
	
	\begin{enumerate}
		\item First item
		\item Second item
		\item Third item
	\end{enumerate}
	
	\begin{theorem}\label{T1}
		There are infinitely many prime numbers.
	\end{theorem}
	
	In Theorem \ref{T1}, we state that there are infinitely many
	prime numbers.
	
	\section{Science!}
	
	Einstein's famous mass-energy equivalence equation is
	\(E = mc^2\). This equation describes the relationship between
	mass and energy \citep{einstein1905}.
	
	\section{Equations and Cross-References}
	
	The time-dependent Schrödinger equation is given by
	Equation \ref{eq:schrodinger}. Schrödinger's 1926 paper
	discusses the foundations of wave mechanics
	\citep{schrodinger1926}.
	
	\begin{equation}
		i\hbar\frac{\partial}{\partial t}\Psi(\mathbf{r},t)
		=
		\left[
		-\frac{\hbar^2}{2\mu}\nabla^2
		+ V(\mathbf{r},t)
		\right]\Psi(\mathbf{r},t)
		\label{eq:schrodinger}
	\end{equation}
	
	\bibliographystyle{agu}
	\bibliography{physics}
	
\end{document}
